\documentclass[tikz,border=5pt]{standalone}
\usepackage{ctex}
\usepackage{amsmath,amssymb}
\usetikzlibrary{arrows.meta,positioning}
% Shared visual language for LFS architecture diagrams.
\RequirePackage{../../mymath}

\definecolor{LFSBlue}{HTML}{2563EB}
\definecolor{LFSBlueSoft}{HTML}{DBEAFE}
\definecolor{LFSGreen}{HTML}{15803D}
\definecolor{LFSGreenSoft}{HTML}{DCFCE7}
\definecolor{LFSOrange}{HTML}{C2410C}
\definecolor{LFSOrangeSoft}{HTML}{FFEDD5}
\definecolor{LFSRed}{HTML}{B91C1C}
\definecolor{LFSRedSoft}{HTML}{FEE2E2}
\definecolor{LFSPurple}{HTML}{7E22CE}
\definecolor{LFSPurpleSoft}{HTML}{F3E8FF}
\definecolor{LFSGray}{HTML}{475569}
\definecolor{LFSGraySoft}{HTML}{F1F5F9}

\tikzset{
  lfs picture/.style={
    font=\small,
    node distance=8mm and 12mm,
    >=Latex,
    every path/.style={line cap=round, line join=round}
  },
  block/.style={
    draw=LFSBlue,
    fill=LFSBlueSoft,
    rounded corners=2pt,
    line width=0.8pt,
    align=center,
    inner xsep=7pt,
    inner ysep=5pt,
    minimum height=10mm
  },
  compute/.style={
    block,
    draw=LFSPurple,
    fill=LFSPurpleSoft
  },
  storage/.style={
    block,
    draw=LFSGreen,
    fill=LFSGreenSoft
  },
  interface/.style={
    block,
    draw=LFSOrange,
    fill=LFSOrangeSoft
  },
  risk/.style={
    block,
    draw=LFSRed,
    fill=LFSRedSoft
  },
  neutral/.style={
    block,
    draw=LFSGray,
    fill=LFSGraySoft
  },
  decision/.style={
    diamond,
    draw=LFSOrange,
    fill=LFSOrangeSoft,
    line width=0.8pt,
    align=center,
    inner sep=2pt,
    aspect=2
  },
  group/.style={
    draw=LFSGray,
    rounded corners=3pt,
    dashed,
    line width=0.7pt,
    inner sep=6pt
  },
  arrow/.style={
    -{Latex[length=2.4mm,width=1.6mm]},
    draw=LFSGray,
    line width=0.9pt
  },
  data arrow/.style={
    arrow,
    draw=LFSBlue
  },
  control arrow/.style={
    arrow,
    draw=LFSOrange,
    dashed
  },
  residual/.style={
    arrow,
    draw=LFSGreen,
    rounded corners=4pt
  },
  label text/.style={
    font=\scriptsize,
    text=LFSGray,
    fill=white,
    inner sep=1.5pt
  },
  title/.style={
    font=\bfseries,
    text=LFSGray,
    align=center
  }
}


\begin{document}
\begin{tikzpicture}[my picture, >=Latex, node distance=10mm]

% Style definitions
\tikzset{
  leaf/.style={draw=MyBlue, fill=MyBlueSoft, rounded corners=2pt, font=\small\ttfamily, inner sep=4pt, minimum width=9mm, minimum height=7mm},
  intermediate/.style={draw=MyPurple, fill=MyPurpleSoft, rounded corners=2pt, font=\small\ttfamily, inner sep=4pt, minimum width=14mm, minimum height=7mm},
  root/.style={draw=MyGreen, fill=MyGreenSoft, rounded corners=3pt, font=\small\bfseries, inner sep=5pt, minimum height=8mm},
  rulebox/.style={draw=MyOrange, fill=MyOrangeSoft, rounded corners=2pt, font=\footnotesize, align=left, inner sep=6pt},
  stage/.style={font=\small\bfseries, text=MyGray}
}

% Left: Merge rule hierarchy table / box
\node[stage] (rule_title) at (-0.5, 3.8) {训练阶段：有序合并规则表};
\node[rulebox, below=3mm of rule_title, text width=52mm] (rules) {
  \textbf{优先级队列}（由训练频数决定）：\\[2mm]
  \textbf{Rank 1}: \texttt{('l', 'o')} $\to$ \texttt{'lo'} \hfill (频数最高)\\
  \textbf{Rank 2}: \texttt{('e', 'r')} $\to$ \texttt{'er'}\\
  \textbf{Rank 3}: \texttt{('lo', 'w')} $\to$ \texttt{'low'}\\
  \textbf{Rank 4}: \texttt{('low', '\_')} $\to$ \texttt{'low\_'}\\
  \textbf{Rank 5}: \texttt{('low\_', 'er')} $\to$ \texttt{'lower'}\\[1.5mm]
  \rule{\linewidth}{0.4pt}\\[1.5mm]
  \textit{编码原则}：按规则优先级由小到大\\
  贪心执行全局非重叠字符对合并。
};

% Right: Hierarchical parse tree
\node[stage] (tree_title) at (6.5, 3.8) {推理阶段：层次词元切分树};

% Leaf tokens
\node[leaf] (c1) at (3.5, -0.2) {'l'};
\node[leaf] (c2) at (4.7, -0.2) {'o'};
\node[leaf] (c3) at (5.9, -0.2) {'w'};
\node[leaf] (c4) at (7.1, -0.2) {'\_'};
\node[leaf] (c5) at (8.3, -0.2) {'e'};
\node[leaf] (c6) at (9.5, -0.2) {'r'};

% First level merges
\node[intermediate] (m1) at (4.1, 0.9) {'lo'};
\draw[->, draw=MyPurple, line width=0.8pt] (c1) -- (m1) node[midway, left, font=\tiny, text=MyOrange] {R1};
\draw[->, draw=MyPurple, line width=0.8pt] (c2) -- (m1);

\node[intermediate] (m2) at (8.9, 0.9) {'er'};
\draw[->, draw=MyPurple, line width=0.8pt] (c5) -- (m2) node[midway, left, font=\tiny, text=MyOrange] {R2};
\draw[->, draw=MyPurple, line width=0.8pt] (c6) -- (m2);

% Second level merge
\node[intermediate] (m3) at (5.0, 1.8) {'low'};
\draw[->, draw=MyPurple, line width=0.8pt] (m1) -- (m3) node[midway, left, font=\tiny, text=MyOrange] {R3};
\draw[->, draw=MyPurple, line width=0.8pt] (c3) -- (m3);

% Third level merge
\node[intermediate] (m4) at (6.0, 2.7) {'low\_'};
\draw[->, draw=MyPurple, line width=0.8pt] (m3) -- (m4) node[midway, left, font=\tiny, text=MyOrange] {R4};
\draw[->, draw=MyPurple, line width=0.8pt] (c4) -- (m4);

% Final merge
\node[root] (m5) at (7.45, 3.6) {词元: \texttt{'lower'}};
\draw[->, draw=MyGreen, line width=1pt] (m4) -- (m5) node[midway, above left, font=\tiny, text=MyOrange] {R5};
\draw[->, draw=MyGreen, line width=1pt] (m2) to[out=90, in=-50] (m5);

% Annotation at bottom
\node[font=\footnotesize, text=MyGray, align=center] at (6.5, -0.9) {
  原始序列由 6 个基础字符组成 $\xrightarrow{\text{5 轮合并}}$ 压缩为 1 个完整词元
};

\end{tikzpicture}
\end{document}
